% !PRINTED 387
\weberpage{387}

%<<CONT
so $b$ has one of the values 1, 2, 3 \ldots $\lambda$, but with the restriction that for even $D$ only the even, and for odd $D$ only the odd among these numbers are to be taken for $b$. Then $a$ and $c$ are to be determined so that

\weq{6}{\frac{D-b^2}{4} = ac,}

and so that

\weq{7}{\frac{\lambda-b+1}{2} \leqq c \leqq \frac{\lambda+b}{2},}

and $a$ must lie between the same limits. Only one of these two limiting values is an integer; the other may be replaced, for the lower limit, by the next greater integer, and for the upper limit, by the next smaller. Thus for $b$ only a finite number of values is admissible; then $(D-b^2):4$ is decomposable into two factors in only a finite number of ways, and of these decompositions only those are to be retained in which both factors satisfy the condition (7). Moreover, such combinations are also to be discarded in which $a$, $b$, $c$ acquire a common divisor. Herein lies the means of actually determining all reduced numbers for a given $D$.

For what follows we shall, for the numbers $\omega$ so determined, use the symbol

\weq{8}{\omega = \frac{\sqrt{D}+b}{2c} = \frac{2a}{\sqrt{D}-b} = \{a, b, c\}}

to denote them.

From 3. \S.\,124 and 4. \S.\,125 there now results the following important theorem:

\gesperrt{5. The continued-fraction expansion of a reduced number $\omega$ is periodic,}

that is, if $\omega$ is expanded in a continued fraction

\weq{9}{\omega = (\alpha_0, \alpha_1, \alpha_2 \ldots)}

then after a certain number of partial denominators the same partial denominators always recur in the same order; thus, if the period consists of $\nu$ terms, then
\[
\begin{array}{ccccc}
\alpha_0 &=& \alpha_\nu &=& \alpha_{2\nu} \ldots\\
\alpha_1 &=& \alpha_{\nu+1} &=& \alpha_{2\nu+1} \ldots\\
\multicolumn{5}{c}{\cdots\cdots\cdots\cdots\cdots\cdots\cdots\cdots\cdots}\\
\alpha_{\nu-1} &=& \alpha_{2\nu-1} &=& \alpha_{3\nu-1} \ldots,
\end{array}
\]
%CONT>>
